\title{【サンプル】AI開発ワークフロー}
\tag{サンプル, 生成AI, 開発プロセス, 自動化}
\repos{https://github.com/nebuka-radix/homepage}
\summary{変更依頼からレビューまでを整理した、開発ワークフローの紹介サンプルです。}

\fig{thumbnail.png}{表示確認用のサンプル模式図。実際のアプリ画面ではありません。}

\begin{portfolio}

\section{このページについて}

この作品は、画面を持つアプリだけでなく、
開発手順や仕組みもポートフォリオとして紹介するためのサンプルです。

効果や削減率などの実績を示すものではありません。
リポジトリへのリンクは、本ホームページのソースコードを指しています。

\section{想定する課題}

生成AIへの依頼内容や変更後の確認方法が属人化すると、
チーム内で作業品質をそろえることが難しくなります。

このサンプルでは、依頼、変更、検証、レビューという
作業の段階を明確にすることを想定しています。

\section{ワークフロー}

\subsection{作業状況の整理}

依頼内容、変更候補、確認結果を分けて管理する構成です。

\fig{thumbnail.png}{作業を段階別に整理したボードの模式図。実際のサービス画面ではありません。}

\subsection{人による確認を含む処理}

生成された変更をそのまま公開するのではなく、
検証とレビューを挟む流れを想定しています。

\fig{flow.png}{依頼、変更、検証、レビューの段階を表した模式図。各段階の詳細は説明文で補います。}

\section{紹介形式の確認}

このサンプルでは、成果物の見た目だけでなく、
課題と進め方を文章と図で説明する形式を確認できます。

\end{portfolio}
